\slajdid{09_1-s023}
\begin{frame}[label=09_1-s023]{Održivost naponskih nivoa}\setlength{\parskip}{2pt}\setlength{\abovedisplayskip}{4pt}\setlength{\belowdisplayskip}{4pt}\setlength{\abovedisplayshortskip}{3pt}\setlength{\belowdisplayshortskip}{3pt}
% element: 09_1-s023:e01
% element: 09_1-s023:e02
% element: 09_1-s023:e03
% element: 09_1-s023:e04
% element: 09_1-s023:e05
\begin{columns}[T,onlytextwidth]\column{.28\textwidth}\centering \begin{tikzpicture}[x=.67cm,y=.67cm,font=\sitneoznake,>=Latex,remember picture,line width=.7pt]
\draw[->](-.13,0)--(3.45,0)node[right]{$V_I$};\draw[->](0,-.13)--(0,2.9)node[above]{$V_O$};\node[left]at(0,2.4){$V_{OH}$};\node[left]at(0,.13){$V_{OL}$};
\draw(0,2.4)--(.28,2.4)..controls(.6,2.4)and(.65,2.05)..(.83,1.65)..controls(1.15,.7)and(1.3,.4)..(1.6,.3)--(2.85,.13);
\draw[dashed](.58,-.1)node[below]{$V_{IL}$}--(.58,2.26);\draw[dashed](1.54,-.1)node[below]{$V_{IH}$}--(1.54,.36);\node[below]at(2.85,0){$V_{DD}$};\draw[dashed](0,.13)--(2.85,.13);
\draw[DEcrvena,dashed](0,2.26)--(.58,2.26);\draw[DEcrvena,dashed](0,.36)--(1.54,.36);
\draw(.4,2.44)--(.78,2.06);\node[right]at(.73,2.43){$a{=}{-}1$};\draw(1.31,.59)--(1.77,.13);\node[right]at(2.5,.85){$a{=}{-}1$};\coordinate(s23from)at(0,.36);
\draw[->,DEcrvena](2.2,1.5)--(1.54,.42);\node[above]at(2.05,1.65){$V_{O(IH)}$};\end{tikzpicture}
\column{.68\textwidth}\textbf{Logička nula mora zakočiti sledeći tranzistor}
\[\frac{V_{DD}}{k_nR_L(V_{DD}-V_{Tn})+1}<V_{Tn}\]
\[\alert{V_{O(IH)}<V_{IL}},\qquad\text{poželjno}\quad\alert{V_{O(IH)}<V_{Tn}}\]
\[\sqrt{\frac{2V_{DD}}{3k_nR_L}}<V_{Tn}\]
\end{columns}\medskip Za konstantno napajanje podešavamo $k_nR_L$, uz $\beta=k_nR_L/V_{DD}$.\[\sqrt{\frac2{3\beta}}<V_{Tn}\ \Rightarrow\ \beta>\frac2{3V_{Tn}^2}\]
Potreban je i uslov \alert{$V_{O(IL)}>V_{IH}$}.
\end{frame}
